\documentclass[11pt]{article}

\usepackage[a4paper,margin=2.7cm]{geometry}
\usepackage{amsmath,amssymb,amsthm,mathtools}
\usepackage{enumitem}
\usepackage{hyperref}

\newtheorem{theorem}{Theorem}[section]
\newtheorem{proposition}[theorem]{Proposition}
\newtheorem{lemma}[theorem]{Lemma}
\newtheorem{corollary}[theorem]{Corollary}
\newtheorem{definition}[theorem]{Definition}
\newtheorem{remark}[theorem]{Remark}

\newcommand{\C}{\mathbb C}
\newcommand{\D}{\mathbb D}
\newcommand{\R}{\mathbb R}
\newcommand{\dd}{\,d}
\newcommand{\pbar}{\partial_{\bar z}}
\newcommand{\p}{\partial}
\newcommand{\1}{\mathbf 1}
\newcommand{\Log}{\operatorname{Log}}

\title{Asymptotic experiments for the analytic two-boundary inverse problem}
\author{}
\date{}

\begin{document}
\maketitle

\section{Purpose}

These notes record what one can rigorously extract from the boundary moment identities in
the inverse Gaussian localization problem, under the following regularity assumption:
\[
  \partial U=\Gamma_{\mathrm{out}}\cup \Gamma_{\mathrm{in}}
\]
has exactly two connected components, and both are real-analytic Jordan curves. No
stronger global regularity or algebraicity of the boundary is assumed.

The goal is not to give a finished rigidity proof. The point is to isolate the correct
asymptotic mechanism. A real-analytic boundary gives a complex collar. It does not, by
itself, give a global meromorphic Schwarz function or a finite singularity structure.
Consequently the correct asymptotics are Darboux/steepest-descent asymptotics controlled
by the first complex obstruction to continuing or deforming the boundary integrals, not by
the largest Euclidean radius attained on the real boundary.

\section{Moment and boundary identities}

For the Gaussian window and the Hermite eigenfunction \(h_k\), the Bargmann-side
eigenfunction identity gives
\begin{equation}\label{eq:bulk-moments}
  \int_U z^n |z|^{2k}e^{-\pi |z|^2}\dd A(z)=0,
  \qquad n\ge 1.
\end{equation}
Define
\begin{equation}\label{eq:Hk-def}
  H_k(s):=s^{-(k+1)}\int_0^s t^k e^{-\pi t}\dd t .
\end{equation}
The singularity at \(s=0\) is removable, and \(H_k\) is entire, since
\[
  H_k(s)=\sum_{m=0}^\infty
  \frac{(-\pi)^m}{m!(m+k+1)}s^m .
\]
Moreover,
\[
  \pbar\!\left(z^{n+k}\bar z^{\,k+1}H_k(|z|^2)\right)
  =
  z^n |z|^{2k}e^{-\pi |z|^2}.
\]
Therefore Cauchy--Green gives, for every \(n\ge 1\),
\begin{equation}\label{eq:boundary-moments}
  \int_{\partial U} z^{n+k}\bar z^{\,k+1}H_k(|z|^2)\dd z=0.
\end{equation}
Writing the two boundary components with the orientation induced by \(U\),
\begin{equation}\label{eq:two-component-split}
  I_n^{\mathrm{out}}+I_n^{\mathrm{in}}=0,
\end{equation}
where
\begin{equation}\label{eq:component-integral}
  I_n^\Gamma
  :=
  \int_\Gamma z^{n+k}\bar z^{\,k+1}H_k(|z|^2)\dd z .
\end{equation}

\section{Why the naive real Laplace asymptotic is false}

Let a component be parametrized by a real-analytic immersion
\[
  \gamma:\R/2\pi\mathbb Z\to \C\setminus\{0\}.
\]
Then
\begin{equation}\label{eq:param-integral}
  I_n^\Gamma=\int_0^{2\pi}\gamma(t)^n A(t)\dd t,
\end{equation}
where
\[
  A(t)=\gamma(t)^k\overline{\gamma(t)}^{\,k+1}
       H_k(|\gamma(t)|^2)\gamma'(t).
\]
It is tempting to set \(R(t)=|\gamma(t)|\) and claim that a nonconstant maximum
\(R(t_0)=R_*\) contributes
\[
  R_*^n n^{-1/(2m)}e^{in\arg\gamma(t_0)}
\]
up to a nonzero constant. This is not correct in general.

Indeed, near \(t_0\) write
\[
  \gamma(t)=R(t)e^{i\varphi(t)} .
\]
If \(R'(t_0)=0\), then regularity of the curve gives
\[
  \gamma'(t_0)=e^{i\varphi(t_0)}R_*\,i\varphi'(t_0).
\]
Since \(\gamma'(t_0)\ne0\), necessarily
\[
  \varphi'(t_0)\ne0.
\]
Thus the local model contains a large linear oscillation:
\[
  \gamma(t)^n
  =
  R_*^n e^{in\varphi(t_0)}
  \exp\{-na(t-t_0)^{2m}+inb(t-t_0)+\cdots\},
  \qquad b\ne0.
\]
Already for \(m=1\),
\[
  \int_{\R}e^{-na t^2+inbt}\dd t
  =
  \sqrt{\frac{\pi}{na}}\,
  e^{-b^2n/(4a)}.
\]
Hence the real maximum can be exponentially smaller than \(R_*^n n^{-1/2}\). The
dominant contribution has moved to a complex saddle. The correct asymptotics must be
complex.

\section{The correct local analytic setup}

Because \(\gamma\) is real analytic, there is a strip
\[
  \mathcal S_\epsilon:=\{t\in\C/2\pi\mathbb Z:|\Im t|<\epsilon\}
\]
to which \(\gamma\) extends holomorphically. The conjugate boundary value is continued by
the reflected analytic function
\[
  \gamma^*(t):=\overline{\gamma(\bar t)} ,
\]
which satisfies \(\gamma^*(t)=\overline{\gamma(t)}\) for real \(t\). Thus the amplitude is
continued as
\[
  A(t):=\gamma(t)^k\gamma^*(t)^{k+1}
       H_k(\gamma(t)\gamma^*(t))\gamma'(t).
\]
Shrinking \(\epsilon\), we may assume \(\gamma(t)\ne0\) and \(A(t)\) is holomorphic in the
strip. Choose a branch of \(\Psi=\Log\gamma\) on the universal cover of this strip. Then
\[
  I_n^\Gamma=\int_{\R/2\pi\mathbb Z} e^{n\Psi(t)}A(t)\dd t .
\]
The possible leading contributions are:
\begin{enumerate}[label=(\roman*)]
\item complex critical points satisfying
\[
  \Psi'(\tau)=0
  \quad\Longleftrightarrow\quad
  \gamma'(\tau)=0;
\]
\item singularities of the analytic continuation of \(\gamma\), \(A\), or the chosen
branch of \(\Psi\);
\item boundary points of the maximal continuation domain, if no explicit singularity
description is available.
\end{enumerate}
Thus real analyticity alone gives a method, but not a numerical list of dominant
contributions. To get a closed formula one must know, or prove, which critical points or
singularities are first encountered by an admissible contour deformation.

\begin{definition}[Admissible deformation and exponential scale]
Fix a holomorphic continuation domain \(X\) for \(\gamma\) and \(A\), containing the real
cycle \(C_0=\R/2\pi\mathbb Z\), and assume \(\gamma\ne0\) on \(X\). An admissible cycle is a
cycle \(C\subset X\), homologous to \(C_0\) in \(X\). Define
\[
  M(C):=\sup_{t\in C}|\gamma(t)|,
  \qquad
  \rho_X(\Gamma):=\inf_{C\sim C_0}M(C).
\]
\end{definition}

If the infimum is achieved on a smooth contour away from singularities and away from
critical points of \(\Log\gamma\), then the contour can be pushed a little farther and the
infimum was not sharp. Therefore a sharp contour is obstructed by a saddle
\(\gamma'(\tau)=0\) or by the boundary of the continuation domain. In this sense the
geometric exponential scale is
\[
  \rho_X(\Gamma)
  =
  \max |\gamma(\tau)|
\]
over the dominant obstruction points \(\tau\) on the optimal deformation. This is the
right replacement for the incorrect quantity \(\max_{\Gamma}|z|\).

\section{Saddle asymptotics}

Here is the precise local statement used once the dominant obstruction is a saddle.

\begin{proposition}[Nondegenerate saddle contribution]\label{prop:saddle}
Let \(X\) be a continuation domain for \(\Psi=\Log\gamma\) and \(A\). Suppose an
admissible deformation of the real cycle can be made into a steepest-descent contour
through finitely many points \(\tau_1,\dots,\tau_N\), and that:
\begin{enumerate}[label=(\roman*)]
\item \(\Psi'(\tau_j)=0\);
\item \(\Psi''(\tau_j)\ne0\);
\item the values \(\Re\Psi(\tau_j)\) are equal and dominate all other portions of the
deformed contour;
\item \(A(\tau_j)\ne0\).
\end{enumerate}
Then
\begin{equation}\label{eq:saddle-asymptotic}
  I_n^\Gamma
  \sim
  \sum_{j=1}^N
  \varepsilon_j
  A(\tau_j)e^{n\Psi(\tau_j)}
  \left(\frac{2\pi}{-n\Psi''(\tau_j)}\right)^{1/2}
  \left(1+\frac{c_{j,1}}{n}+\frac{c_{j,2}}{n^2}+\cdots\right),
\end{equation}
where the signs/phases \(\varepsilon_j\) and square-root branches are determined by the
orientation of the steepest-descent contour.
\end{proposition}

\begin{remark}
In geometric terms, the exponential scale in \eqref{eq:saddle-asymptotic} is
\[
  \rho_\Gamma=\max_j |\gamma(\tau_j)|.
\]
If several saddles have the same modulus, the leading term is a finite exponential sum
\[
  \rho_\Gamma^n n^{-1/2}
  \sum_j c_j e^{in\arg \gamma(\tau_j)} .
\]
Degenerate saddles give the same structure with \(n^{-1/2}\) replaced by a fractional
power \(n^{-1/p}\), where \(p\) is the first nonzero order of \(\Psi-\Psi(\tau_j)\).
\end{remark}

\section{Schwarz-function/Darboux formulation}

There is an equivalent formulation directly in the \(z\)-plane. Since \(\Gamma\) is
real analytic, it has a Schwarz function \(S_\Gamma\) in a collar:
\[
  S_\Gamma(z)=\bar z,\qquad z\in\Gamma.
\]
In that collar,
\begin{equation}\label{eq:B-schwarz}
  I_n^\Gamma=\int_\Gamma z^n B_\Gamma(z)\dd z,
  \qquad
  B_\Gamma(z):=z^k S_\Gamma(z)^{k+1}H_k(zS_\Gamma(z)).
\end{equation}

Assume \(B_\Gamma\) can be analytically continued to the region swept when \(\Gamma\) is
deformed to another curve \(\Gamma'\), except for isolated singularities
\(\alpha_1,\dots,\alpha_N\). Then Cauchy's theorem gives the exact identity
\begin{equation}\label{eq:darboux-exact}
  I_n^\Gamma
  =
  2\pi i\sum_{j=1}^N \operatorname{Res}_{z=\alpha_j}
  \left(z^nB_\Gamma(z)\right)
  +
  \int_{\Gamma'}z^nB_\Gamma(z)\dd z .
\end{equation}
If \(\sup_{\Gamma'}|z|<\rho\), where
\[
  \rho:=\max_j|\alpha_j|,
\]
then the last term is \(O((\rho-\eta)^n)\) for some \(\eta>0\), and the asymptotics are
given by the local singular expansions of \(B_\Gamma\) at the singularities with
\(|\alpha_j|=\rho\).

\begin{proposition}[Darboux forms]\label{prop:darboux-forms}
Suppose \(\alpha\ne0\) is a dominant singularity of \(B_\Gamma\).
\begin{enumerate}[label=(\roman*)]
\item If \(B_\Gamma\) has a pole of order \(q\) at \(\alpha\), then its contribution has
the form
\[
  \alpha^n P_\alpha(n),
\]
where \(P_\alpha\) is a polynomial of degree at most \(q-1\).
\item If
\[
  B_\Gamma(z)
  =
  C_\alpha (1-z/\alpha)^\beta
  +\text{less singular terms},
  \qquad \beta\notin\mathbb Z_{\ge0},
\]
then its contribution is
\[
  C_\alpha'
  \alpha^n n^{-\beta-1}
  \left(1+O(n^{-1})\right),
\]
with the constant determined by the local branch.
\item If the singularity is essential, the local saddle problem near \(\alpha\) must be
solved. A standard normal form is
\[
  B_\Gamma(z)\approx C(z-\alpha)^{-q}
  \exp\!\left(\frac{\beta}{z-\alpha}\right),
\]
which leads to saddles satisfying
\[
  \frac{n}{\alpha}+O(n(z-\alpha))-\frac{\beta}{(z-\alpha)^2}=0
\]
and therefore to contributions of the type
\[
  \alpha^n
  \exp\!\left(\pm 2\sqrt{\frac{\beta n}{\alpha}}\right)
  n^{q/2-3/4}
  \times(\text{explicit phase and constant}),
\]
provided the relevant steepest paths are accessible. The exact sign and branch are part of
the local steepest-descent computation.
\end{enumerate}
\end{proposition}

\begin{remark}
The essential case is not artificial here. Even if \(S_\Gamma\) has only a pole, the
factor \(H_k(zS_\Gamma(z))\) is entire but not polynomial, so composing it with a pole of
\(S_\Gamma\) can create an essential singularity. In favorable sectors one may replace
\(H_k(s)\) by its large-\(s\) asymptotic, but this is a sectorial statement and must be
checked against the contour geometry.
\end{remark}

\section{What step 5 should actually prove}

For two components, the boundary identity is
\[
  I_n^{\mathrm{out}}+I_n^{\mathrm{in}}=0.
\]
The asymptotic strategy can prove that the outer boundary is a centered circle if it
establishes the following structure.

\begin{theorem}[Conditional separation criterion]\label{thm:conditional-separation}
Assume both boundary components are real analytic. Suppose that, for each component, one
has chosen a continuation/deformation domain giving true asymptotics
\begin{align}
  I_n^{\mathrm{out}}
  &=
  \rho_{\mathrm{out}}^n n^{a_{\mathrm{out}}}
  \left(\sum_{\ell=1}^{L}c_\ell e^{in\theta_\ell}+o(1)\right),
  \label{eq:out-asymp}\\
  I_n^{\mathrm{in}}
  &=
  O\!\left(\rho_{\mathrm{in}}^n n^{a_{\mathrm{in}}}\right),
  \label{eq:in-asymp}
\end{align}
with \(c_\ell\ne0\) and distinct \(e^{i\theta_\ell}\). If
\[
  \rho_{\mathrm{in}}<\rho_{\mathrm{out}},
\]
then \eqref{eq:two-component-split} is impossible. Hence the hypothesis that produced the
nonzero outer asymptotic must fail.
\end{theorem}

\begin{proof}
Divide \eqref{eq:two-component-split} by
\(\rho_{\mathrm{out}}^n n^{a_{\mathrm{out}}}\). The inner contribution tends to \(0\).
Thus
\[
  \sum_{\ell=1}^{L}c_\ell e^{in\theta_\ell}\to0.
\]
By Cesaro averaging after multiplication by \(e^{-in\theta_j}\), each coefficient
\(c_j\) must vanish, a contradiction.
\end{proof}

The geometric content is therefore the strict separation
\[
  \rho_{\mathrm{in}}<\rho_{\mathrm{out}}.
\]
Here \(\rho_{\mathrm{out}}\) is not \(\max_{\Gamma_{\mathrm{out}}}|z|\). It is the modulus
of the dominant obstruction for the analytic continuation of the outer component integral:
either \(|\gamma_{\mathrm{out}}(\tau)|\) at a dominant complex saddle
\(\gamma'_{\mathrm{out}}(\tau)=0\), or \(|\alpha|\) at a dominant Schwarz-function
singularity \(z=\alpha\). The same definition applies to \(\rho_{\mathrm{in}}\).

\section{What real analyticity alone gives, and what it does not give}

Real analyticity gives:
\begin{enumerate}[label=(\roman*)]
\item a nontrivial complex collar for each boundary component;
\item local Schwarz functions \(S_{\mathrm{out}}\) and \(S_{\mathrm{in}}\);
\item a legitimate steepest-descent/Darboux problem for each component integral;
\item true asymptotics once the relevant saddles or singularities in the chosen
continuation domain are identified.
\end{enumerate}

Real analyticity alone does not give:
\begin{enumerate}[label=(\roman*)]
\item a global meromorphic Schwarz function;
\item a finite list of singularities;
\item equality between the asymptotic scale and the real boundary radius
\(\max_\Gamma |z|\);
\item the separation \(\rho_{\mathrm{in}}<\rho_{\mathrm{out}}\).
\end{enumerate}

Thus a rigorous proof should not state that the outer component has leading size
\[
  R_*^n n^{-1/(2m)}
  \quad\text{with } R_*=\max_{\Gamma_{\mathrm{out}}}|z|.
\]
The right statement is conditional on the complexified geometry:
\[
  I_n^\Gamma
  \text{ is controlled by dominant accessible complex saddles or singularities.}
\]

\section{Suggested proof program}

The most plausible rigorous route is:
\begin{enumerate}[label=\arabic*.]
\item Derive \eqref{eq:boundary-moments} from the Hermite eigenfunction identity.
\item For each analytic component, introduce its local Schwarz function and write
\[
  I_n^\Gamma=\int_\Gamma z^nB_\Gamma(z)\dd z.
\]
\item Determine a continuation/deformation domain on the side relevant to the annular
region \(U\). This is an actual geometric input, not a formal consequence of real
analyticity.
\item Apply Darboux if the first obstructions are Schwarz singularities; apply complex
steepest descent if the first obstructions are saddles of \(\Log\gamma\).
\item Prove the strict separation
\[
  \rho_{\mathrm{in}}<\rho_{\mathrm{out}}.
\]
\item Use Theorem~\ref{thm:conditional-separation} to force the outer boundary not to have
any nonradial dominant obstruction. If this can be made for every possible obstruction,
the outer boundary must be a circle centered at the origin.
\item Once the outer boundary is a centered circle, the angular-derivative propagation
argument becomes robust: in a circular collar, \(w=\partial_\theta u\) is harmonic and has
both \(w=0\) and \(\partial_r w=0\) on the outer circle, hence \(w\equiv0\) in the collar
and then in the connected domain by unique continuation. This yields radiality.
\end{enumerate}

\section{Bottom line}

The correct asymptotics are available under real-analytic boundary regularity, but their
exponential scales are complex-geometric:
\[
  \rho_\Gamma
  =
  \text{modulus of the dominant accessible saddle/singularity of the component integral}.
\]
They are not determined by \(\max_\Gamma |z|\). A complete proof of radiality by this
method therefore needs a new geometric theorem comparing the dominant complex
obstructions of the outer and inner analytic boundary components.

\section{Why weighted quadrature does not automatically give a meromorphic Schwarz function}

One might hope to strengthen the previous program as follows. Compact support gives a
weighted quadrature identity; perhaps this forces the ordinary Schwarz function to extend
meromorphically, after which the asymptotics reduce to residues. This is too optimistic in
general.

Suppose
\[
  \Delta u=f\,\1_U-c\delta_0
\]
has a compactly supported solution, with \(f\) real analytic near \(\partial U\). Testing
against holomorphic polynomials gives
\begin{equation}\label{eq:weighted-quadrature}
  \int_U p(z)f(z,\bar z)\dd A(z)=c\,p(0).
\end{equation}
Equivalently, for \(\zeta\) outside \(U\),
\begin{equation}\label{eq:weighted-cauchy}
  C_{f,U}(\zeta):=\int_U\frac{f(z,\bar z)}{\zeta-z}\dd A(z)
  =
  \frac{c}{\zeta}.
\end{equation}
Thus the data control the weighted Cauchy transform \(C_{f,U}\), not the ordinary Cauchy
transform
\[
  C_U(\zeta):=\int_U\frac{\dd A(z)}{\zeta-z}.
\]
Classical quadrature-domain theory gives meromorphic Schwarz functions from finite
quadrature identities for \(dA\), or equivalently from rationality/finite-rank structure of
\(C_U\). The identity \eqref{eq:weighted-cauchy} is a different statement.

Let \(F(z,\bar z)\) be a local analytic \(\pbar\)-primitive of \(f\):
\[
  \pbar F(z,\bar z)=f(z,\bar z).
\]
On an analytic boundary component, write \(S(z)=\bar z\). Then Cauchy--Green gives
\[
  \int_U z^n f(z,\bar z)\dd A(z)
  =
  \frac{1}{2i}\int_{\partial U}z^nF(z,\bar z)\dd z
  =
  \frac{1}{2i}\int_{\partial U}z^nF(z,S(z))\dd z.
\]
The boundary object controlled by the weighted moment identities is therefore
\begin{equation}\label{eq:G-weighted-schwarz}
  G(z):=F(z,S(z)),
\end{equation}
not \(S(z)\) itself.

The missing implication would be
\begin{equation}\label{eq:bad-step2}
  G(z)=F(z,S(z))\ \text{meromorphic}
  \quad\Longrightarrow\quad
  S(z)\ \text{meromorphic}.
\end{equation}
This implication is not formal. Its truth depends on the dependence of \(F\) on the second
variable.

\begin{enumerate}[label=(\roman*)]
\item If \(f=1\), we may take \(F(z,\eta)=\eta\). Then \(G=S\), and the weighted identity is
the usual quadrature identity. This is the classical case in which meromorphic
continuation of the Schwarz function is natural.

\item More generally, if
\[
  F(z,\eta)=a(z)\eta+b(z),
\]
with \(a\) holomorphic and nonvanishing in the continuation region, then
\[
  S(z)=\frac{G(z)-b(z)}{a(z)}.
\]
In this affine case, meromorphicity of \(G\) implies meromorphicity of \(S\).

\item If \(F\) is polynomial in \(\eta\) of degree \(d>1\), then \(G=F(z,S)\) gives an
algebraic equation for \(S\), not a meromorphic representation. One may at best expect
\(S\) to be algebraic over the meromorphic field generated by \(G\), and branching can
appear.

\item If \(F\) is non-polynomial analytic in \(\eta\), then \eqref{eq:bad-step2} is
generally the wrong target. Inverting \(F(z,\cdot)\) can produce multivalued branches and
essential complications even when \(G\) is meromorphic.
\end{enumerate}

For the Hermite/Gaussian boundary identity, the natural object is not \(S\) but
\begin{equation}\label{eq:weighted-hermite-object}
  B(z)=z^kS(z)^{k+1}H_k(zS(z)).
\end{equation}
The function \(H_k\) is entire and non-polynomial. Hence even if \(S\) had only a pole,
the composition \(H_k(zS(z))\) could produce an essential singularity. Conversely, knowing
that \(B\) has controlled continuation does not by itself make \(S\) meromorphic.

Thus the realistic replacement for the failed Step 2 is:
\[
  \text{prove controlled continuation/asymptotics for }F(z,S(z))
  \text{ or }B(z),
\]
rather than trying to prove meromorphic continuation of the ordinary Schwarz function.
Meromorphic \(S\) is plausible only under additional finite-quadrature structure, such as
constant density or affine dependence of a \(\pbar\)-primitive on \(\bar z\).

\section{The pure Gaussian weight}

Now specialize to
\[
  f(z,\bar z)=e^{-\pi |z|^2}=e^{-\pi z\bar z}.
\]
Treating \(z\) and \(\eta=\bar z\) as independent complex variables, a removable
\(\pbar\)-primitive is
\begin{equation}\label{eq:gaussian-primitive}
  F(z,\eta):=
  \begin{cases}
  \displaystyle \frac{1-e^{-\pi z\eta}}{\pi z}, & z\ne0,\\[1.1em]
  \eta, & z=0.
  \end{cases}
\end{equation}
Indeed, the power-series expansion
\[
  \frac{1-e^{-\pi z\eta}}{\pi z}
  =
  \sum_{m=1}^\infty \frac{(-1)^{m+1}\pi^{m-1}}{m!}z^{m-1}\eta^m
\]
shows removability at \(z=0\), and direct differentiation gives
\[
  \partial_\eta F(z,\eta)=e^{-\pi z\eta}.
\]

On an analytic boundary component with Schwarz function \(S\), the controlled boundary
object is therefore
\begin{equation}\label{eq:gaussian-G}
  G(z):=F(z,S(z))
  =
  \frac{1-e^{-\pi zS(z)}}{\pi z},
\end{equation}
again with the removable interpretation at \(z=0\). The moment identity becomes
\begin{equation}\label{eq:gaussian-boundary-moment}
  \int_{\partial U} z^n G(z)\dd z=0,\qquad n\ge1.
\end{equation}
Equivalently, since \(\int_{\partial U}z^{n-1}\dd z=0\) for \(n\ge1\),
\begin{equation}\label{eq:gaussian-exp-boundary}
  \int_{\partial U}z^{n-1}e^{-\pi zS(z)}\dd z=0,\qquad n\ge1.
\end{equation}
This is exactly the \(k=0\) instance of the Hermite boundary identity, because
\[
  H_0(s)=\frac{1-e^{-\pi s}}{\pi s}
  \quad\text{and}\quad
  S(z)H_0(zS(z))
  =
  \frac{1-e^{-\pi zS(z)}}{\pi z}.
\]

This specialization makes the failed meromorphic-Schwarz step especially transparent.
Suppose, optimistically, that the moment identity allowed one to continue \(G\)
meromorphically. Then
\begin{equation}\label{eq:gaussian-M}
  M(z):=e^{-\pi zS(z)}=1-\pi zG(z)
\end{equation}
would be meromorphic. But recovering \(S\) requires taking a logarithm:
\begin{equation}\label{eq:S-from-M}
  S(z)=-\frac{1}{\pi z}\Log M(z).
\end{equation}
This is not a meromorphic operation. Zeros and poles of \(M\) create logarithmic
singularities of \(S\), and nontrivial monodromy may appear. Thus meromorphicity of the
Gaussian weighted object \(G\), or of \(M\), does not imply meromorphicity of the ordinary
Schwarz function.

Conversely, meromorphicity of \(S\) is not naturally compatible with the Gaussian object.
If \(S\) has a pole at \(\alpha\ne0\), then \(zS(z)\) has a pole at \(\alpha\), and
\[
  e^{-\pi zS(z)}
\]
has an essential singularity at \(\alpha\). Hence \(G\) has an essential singularity there,
not a pole. Therefore, for the Gaussian weight, ordinary meromorphic Schwarz continuation
is not the right expected finite-complexity statement. The correct object for asymptotics
is the exponential weighted Schwarz object \(M=e^{-\pi zS}\), or equivalently
\(G=(1-M)/(\pi z)\).

If a continuation of \(S\) has a pole of order \(q\) at a nonzero point \(\alpha\),
\[
  S(z)=a_{-q}(z-\alpha)^{-q}+a_{-q+1}(z-\alpha)^{-q+1}+\cdots,
  \qquad a_{-q}\ne0,
\]
then the Gaussian boundary integrand has local exponential type
\[
  z^nG(z)\dd z
  \sim
  -\frac{z^{n-1}}{\pi}
  \exp\!\left(
    -\pi\alpha a_{-q}(z-\alpha)^{-q}
  \right)\dd z
\]
up to lower singular orders in the exponential. The local saddle equation has leading
form
\begin{equation}\label{eq:gaussian-pole-saddles}
  \frac{n}{\alpha}
  +
  \pi q\alpha a_{-q}(z-\alpha)^{-q-1}
  \approx0,
\end{equation}
so the relevant saddles lie at distance
\begin{equation}\label{eq:gaussian-pole-scale}
  z-\alpha
  \asymp
  n^{-1/(q+1)}.
\end{equation}
Substitution gives contributions of the qualitative form
\begin{equation}\label{eq:gaussian-pole-asymp}
  \alpha^n
  \exp\!\left(C\,n^{q/(q+1)}\right)
  n^\beta
  \left(c_0+c_1 n^{-1/(q+1)}+\cdots\right),
\end{equation}
where \(C,\beta,c_0,\dots\) depend on the chosen saddle, branch, and steepest path. This
is not a residue asymptotic. It is an essential-singularity saddle asymptotic.

\begin{remark}[Two checks]
First, \eqref{eq:gaussian-primitive} is the correct primitive: differentiating the
removable power series in \(\eta\) recovers \(e^{-\pi z\eta}\), and therefore
Cauchy--Green gives \eqref{eq:gaussian-boundary-moment}. Second, the inversion obstruction
is structural: from \(M=e^{-\pi zS}\) one obtains \(S=-(\pi z)^{-1}\Log M\), and logarithms
of meromorphic functions are not meromorphic except in special single-valued nonvanishing
cases. Thus the Gaussian weight does not repair Step 2; it confirms that the asymptotic
object should be weighted/exponential rather than the ordinary Schwarz function.
\end{remark}

\section{Topology: what is independent of simple connectivity}

The continuation of the Gaussian weighted object is local and does not require simple
connectivity. Let
\[
  A(z):=4\partial_z u(z)-F(z,\bar z),
  \qquad
  F(z,\eta)=\frac{1-e^{-\pi z\eta}}{\pi z},
\]
with the removable value \(F(0,\eta)=\eta\). In \(U\setminus\{0\}\),
\[
  \pbar A=4\partial_z\pbar u-\pbar F
  =
  \Delta u-e^{-\pi z\bar z}
  =
  0.
\]
Thus \(A\) is holomorphic on every connected component of \(U\setminus\{0\}\). Now fix an
analytic boundary component \(\Gamma\). The boundary identity below requires one extra
componentwise hypothesis: the complementary component adjacent to \(\Gamma\) must be a
zero region for \(u\). Under this hypothesis, \(C^1\) regularity across \(\Gamma\) gives
\[
  \partial_z u=0\qquad\text{on }\Gamma.
\]
Consequently,
\[
  A(z)=-F(z,S(z))
  =
  \frac{e^{-\pi zS(z)}-1}{\pi z}
  \qquad (z\in\Gamma),
\]
and hence
\begin{equation}\label{eq:M-pde-continuation}
  M(z):=e^{-\pi zS(z)}=1+\pi zA(z)
  \qquad (z\in\Gamma).
\end{equation}
The right-hand side gives a holomorphic continuation of the \(M\)-branch attached to
\(\Gamma\) to the \(U\)-side of that boundary component. This construction uses only the
local PDE, the local analytic Schwarz function near the boundary, and the zero-side
condition for \(u\); it does not use simple connectivity.

This point is important in a multiply connected domain. Compact support forces \(u\equiv0\)
in the unbounded complementary component, by harmonic unique continuation from the open
set where the compactly supported function vanishes. Hence the argument applies to the
outer boundary component. It does not automatically apply to an inner boundary component:
in a bounded hole, \(u\) is harmonic, but it need not vanish identically. On such a
component one only has the identity
\[
  M(z)=1+\pi zA(z)-4\pi z\,\partial_z u(z)
  \qquad (z\in\Gamma),
\]
and the last term is not determined by real analyticity of \(\Gamma\). Thus the
\(M\)-branch continued from the outer boundary need not agree with the local
Schwarz-defined \(M\)-branch on an inner boundary unless an additional zero-side, or
equivalent boundary, condition is proved there.

If a connected component of \(U\) contains the origin, then the distributional statement on
the whole component has to include the point mass. In that case
\[
  \pbar A=-c\delta_0,
\]
and the corrected holomorphic field is
\[
  A_c(z):=A(z)+\frac{c}{\pi z}.
\]
Equivalently,
\[
  A_c=4\partial_z u-F(z,\bar z)+\frac{c}{\pi z}
\]
is holomorphic across \(0\) when the logarithmic singular part has the prescribed strength.
For boundary components away from \(0\), this changes the boundary formula to
\begin{equation}\label{eq:M-pde-continuation-puncture}
  M(z)=1+\pi zA_c(z)-c.
\end{equation}
In an annular component not containing \(0\), no correction is present and
\eqref{eq:M-pde-continuation} is the relevant formula.

What does depend on topology is the attempt to recover the ordinary Schwarz function from
\(M\). Formally,
\[
  S(z)=-\frac{1}{\pi z}\Log M(z).
\]
For this to define a single-valued continuation of \(S\), one needs a single-valued branch
of \(\Log M\). On a multiply connected region this requires not only that \(M\) be
nonzero, but also vanishing of the logarithmic periods:
\begin{equation}\label{eq:log-period-condition}
  \frac{1}{2\pi i}\int_\gamma \frac{M'(z)}{M(z)}\dd z=0
  \qquad
  \text{for every closed cycle }\gamma .
\end{equation}
If these periods do not vanish, \(M\) may be a perfectly good holomorphic continuation of
the weighted object while \(S\) is only a multivalued logarithmic lift.

Topology also enters the contour deformation problem. The asymptotic integrals are
\[
  J_n^\Gamma=\int_\Gamma z^{n-1}M(z)\dd z.
\]
If \(M\) is holomorphic in an annular component of \(U\) between two analytic boundary
curves, then \(z^{n-1}M(z)\dd z\) is a holomorphic one-form there, away from the possible
puncture \(0\). The outer and inner boundary cycles are homologous with opposite
orientations, so
\[
  J_n^{\mathrm{out}}+J_n^{\mathrm{in}}=0
\]
is simply Cauchy's theorem in that annulus. In a multiply connected region one may deform a
contour only within its homology class, not necessarily collapse it to nothing.

If \(0\) lies in the component under consideration, then the point mass contributes the
logarithmic singular part of \(u\), and \(A\) may have a pole at \(0\). This is a puncture
issue, not a simple-connectivity issue. The singular part must be computed before applying
Cauchy's theorem to \(z^{n-1}M(z)\dd z\).

Thus the useful summary is:
\begin{itemize}
\item holomorphic continuation of the \(M\)-branch from a boundary component is local and
topology-free once the adjacent zero-side condition \(u\equiv0\) is available;
\item compact support gives this condition on the unbounded exterior boundary, but not on
inner boundaries surrounding holes;
\item single-valued recovery of \(S\) requires monodromy control, and global contour
deformation requires homology control.
\end{itemize}

\section{Parametrized component asymptotics}

This section refines the older real-maximum argument in the sibling drafts. Those drafts
compare the outer contribution to inner contributions using \(\max_\Gamma |z|\). The
preceding sections show why that comparison is not reliable: the weighted object \(M\)
continues through the annular region, and any asymptotic formula must be tied to genuine
obstructions on the analytic continuation surface, not merely to a chosen real
parametrization.

Let the two analytic boundary components be parametrized by real-analytic immersions
\[
  \gamma_{\mathrm{out}},\gamma_{\mathrm{in}}:\R/2\pi\mathbb Z\to\C\setminus\{0\},
\]
with the orientations induced by \(U\). For each component define the reflected analytic
continuation
\[
  \gamma_j^*(t):=\overline{\gamma_j(\bar t)} ,
  \qquad j\in\{\mathrm{out},\mathrm{in}\}.
\]
On the real axis, \(\gamma_j^*(t)=\overline{\gamma_j(t)}\), and therefore the Gaussian
boundary identity can be written componentwise as
\begin{equation}\label{eq:param-J}
  J_n^{(j)}
  :=
  \int_0^{2\pi}
  \gamma_j(t)^{\,n-1}
  e^{-\pi\gamma_j(t)\gamma_j^*(t)}
  \gamma_j'(t)\dd t,
  \qquad
  J_n^{\mathrm{out}}+J_n^{\mathrm{in}}=0.
\end{equation}
Equivalently, after choosing a branch of
\[
  \Phi_j(t):=\Log\gamma_j(t)
\]
on the relevant continuation strip,
\[
  J_n^{(j)}
  =
  \int e^{n\Phi_j(t)}a_j(t)\dd t,
  \qquad
  a_j(t):=
  e^{-\pi\gamma_j(t)\gamma_j^*(t)}
  \frac{\gamma_j'(t)}{\gamma_j(t)}.
\]
This is the exact parametrized oscillatory integral. Its true asymptotics are obtained by
deforming the \(t\)-contour in the complexified parameter cylinder. The dominant points
are not real maxima of \(|\gamma_j|\). They are accessible singularities of the analytic
data \((\gamma_j,\gamma_j^*)\), or accessible complex saddles of \(\Phi_j\), i.e.
\[
  \Phi_j'(\tau)=0
  \quad\Longleftrightarrow\quad
  \gamma_j'(\tau)=0.
\]
However, such a critical point is useful only if it is a genuine obstruction for the chosen
homology class of contour deformation. If the same contour can be pushed farther in the
\(z\)-plane because \(M\) is holomorphic there, the parametric saddle is a coordinate
artifact and cannot be used as an intrinsic leading term.


There is a useful normalization at such saddles. Let
\[
  E_j(t):=e^{-\pi\gamma_j(t)\gamma_j^*(t)}.
\]
Since
\[
  \gamma_j(t)^{\,n-1}\gamma_j'(t)\dd t
  =
  \frac1n\dd\left(\gamma_j(t)^n\right),
\]
periodicity gives
\begin{equation}\label{eq:param-integration-by-parts}
  J_n^{(j)}
  =
  -\frac1n\int e^{n\Phi_j(t)}E_j'(t)\dd t.
\end{equation}
This identity is important: the naive saddle contribution of \(a_j=E_j\Phi_j'\) cancels at
leading order. Formula \eqref{eq:param-integration-by-parts} is the safer starting point
for saddle asymptotics.

\begin{proposition}[Parametric saddle contribution]\label{prop:param-saddle}
Assume a contour for \(J_n^{(j)}\) can be deformed to a steepest-descent contour through
finitely many dominant saddle points \(\tau_\ell\). Suppose that, at such a point,
\[
  \Phi_j(t)=\Phi_j(\tau_\ell)
  +\lambda_\ell (t-\tau_\ell)^m+O((t-\tau_\ell)^{m+1}),
  \qquad
  m\ge2,\quad \lambda_\ell\ne0,
\]
and that \(E_j'(\tau_\ell)\ne0\). Then the contribution of \(\tau_\ell\) has the form
\begin{equation}\label{eq:param-saddle-form}
  J_{n,\tau_\ell}^{(j)}
  \sim
  e^{n\Phi_j(\tau_\ell)}
  n^{-1-1/m}
  \left(
    C_{\ell,0}
    +C_{\ell,1}n^{-1/m}
    +C_{\ell,2}n^{-2/m}
    +\cdots
  \right).
\end{equation}
The leading constant is the usual one-dimensional steepest-descent constant for the local
normal form
\[
  -\frac1n\int e^{n\lambda_\ell (t-\tau_\ell)^m}E_j'(\tau_\ell)\dd t;
\]
explicitly, it is a nonzero multiple of
\[
  -E_j'(\tau_\ell)\,\Gamma(1/m)\,\lambda_\ell^{-1/m},
\]
with the root and phase determined by the two steepest rays used by the deformed contour.
If \(E_j'\) vanishes to order \(r\) at \(\tau_\ell\), the power \(n^{-1-1/m}\) is replaced
by \(n^{-1-(r+1)/m}\).
\end{proposition}

Thus a finite set of dominant saddles with equal modulus gives
\begin{equation}\label{eq:param-saddle-sum}
  J_n^{(j)}
  \sim
  \rho_j^n n^{-1-1/m_j}
  \sum_{\ell:\,|\gamma_j(\tau_\ell)|=\rho_j}
  C_{\ell,0}e^{in\arg\gamma_j(\tau_\ell)}
\end{equation}
in the nondegenerate equal-type case, with the evident modifications for different saddle
orders or vanishing orders of \(E_j'\). Here
\[
  \rho_j=|\gamma_j(\tau_\ell)|
\]
on the dominant saddle set. This is the precise replacement for the incorrect real
quantity \(\max_{t\in\R}|\gamma_j(t)|\).

If the first obstruction is not a saddle but a singularity of the continued weighted
object, the \(z\)-plane version is cleaner. Let
\[
  M_j(z)=e^{-\pi zS_j(z)}.
\]
If deforming the \(j\)-th contour crosses finitely many pole singularities \(\alpha\) of
\(M_j\), and near such a point
\[
  M_j(z)=\sum_{\ell=1}^{q_\alpha}
  b_{\alpha,-\ell}(z-\alpha)^{-\ell}+O(1),
\]
then the crossed contribution is exactly
\begin{equation}\label{eq:pole-exact}
  2\pi i
  \sum_{\ell=1}^{q_\alpha}
  b_{\alpha,-\ell}
  \binom{n-1}{\ell-1}
  \alpha^{\,n-\ell}.
\end{equation}
Dominant poles are those maximizing \(|\alpha|\) among the crossed singularities.

For algebraic branch singularities,
\[
  M_j(z)\sim b_\alpha(1-z/\alpha)^\beta,
  \qquad \beta\notin\mathbb Z_{\ge0},
\]
the standard Darboux transfer gives contributions of the form
\begin{equation}\label{eq:branch-form}
  C_\alpha\,\alpha^n n^{-\beta-1}
  \left(1+c_1n^{-1}+c_2n^{-2}+\cdots\right),
\end{equation}
with the constant determined by the jump of the chosen branch across the local cut.

For the Gaussian weight, an important singularity model is a pole of the Schwarz function.
If
\[
  S_j(z)=a_{-q}(z-\alpha)^{-q}+\cdots,
  \qquad a_{-q}\ne0,
\]
then \(M_j=e^{-\pi zS_j}\) has an essential singularity at \(\alpha\). The local saddle
calculation can be made more explicit in the model retaining only the leading singular
term. Put \(s=z-\alpha\). The phase is
\[
  \Theta_n(s):=(n-1)\Log(\alpha+s)-\pi(\alpha+s)S_j(\alpha+s).
\]
To leading order,
\[
  \Theta_n'(s)
  =
  \frac{n}{\alpha}
  +
  \pi q\alpha a_{-q}s^{-q-1}
  +\text{lower order terms}.
\]
Thus the \(q+1\) leading saddle locations satisfy
\begin{equation}\label{eq:essential-saddle-locations}
  s_\kappa
  =
  \omega_\kappa
  \left(-\frac{\pi q\alpha^2a_{-q}}{n}\right)^{1/(q+1)}
  \left(1+o(1)\right),
  \qquad \omega_\kappa^{q+1}=1,
\end{equation}
subject to the usual accessibility condition for the deformed contour. In particular the
local saddle scale is
\[
  z-\alpha\asymp n^{-1/(q+1)},
\]
and, for each accessible saddle, the exponent has the refined leading form
\begin{equation}\label{eq:essential-exponent}
  \Theta_n(s_\kappa)
  =
  n\Log\alpha
  +
  C_\kappa n^{q/(q+1)}
  +O(n^{(q-1)/(q+1)}),
\end{equation}
where
\[
  C_\kappa
  =
  \frac{q+1}{q\alpha}\,
  \omega_\kappa
  \left(-\pi q\alpha^2a_{-q}\right)^{1/(q+1)}.
\]
The second derivative at the saddle is of size
\[
  \Theta_n''(s_\kappa)\asymp n^{(q+2)/(q+1)}.
\]
Therefore the corresponding contribution has the form
\begin{equation}\label{eq:essential-form-param}
  \alpha^n
  \exp\!\left(C_\kappa n^{q/(q+1)}\right)
  n^{-\frac{q+2}{2(q+1)}}
  \left(c_0+c_1n^{-1/(q+1)}+c_2n^{-2/(q+1)}+\cdots\right),
\end{equation}
provided the saddle is nondegenerate and no other local terms change the leading order.
The constants \(c_0,c_1,\dots\) and the chosen branch of \(C_\kappa\) are determined by the
accessible steepest path. For \(q=1\), the power is \(n^{-3/4}\), matching the standard
simple-pole/essential-singularity model.

Consequently, the two-component comparison has the following rigorous shape. Once the
dominant accessible saddle/singularity set has been identified for each parametrization,
\[
  J_n^{(j)}
  \sim
  \rho_j^n
  \exp\!\left(C_j n^{\alpha_j}\right)
  n^{\beta_j}
  \sum_\ell c_{j,\ell}e^{in\theta_{j,\ell}}
\]
with \(\alpha_j=0\) in the purely algebraic Darboux/saddle cases and
\(0<\alpha_j<1\) in essential Gaussian cases. The moment identity can force a
contradiction only after one proves a strict dominance statement for this full scale,
lexicographically:
\[
  (\rho_{\mathrm{out}},\Re C_{\mathrm{out}},\beta_{\mathrm{out}})
  >
  (\rho_{\mathrm{in}},\Re C_{\mathrm{in}},\beta_{\mathrm{in}}),
\]
unless the outer component is already a centered circle. This dominance statement is the
geometric input not supplied by real analyticity alone.

\begin{remark}[Scope warning from the sibling files]
The file \texttt{h\_0-loc.tex} records a different, first-order reduction for the
\((h_0,h_0,U)\) localization problem. If that boundary-value interpretation is the one
being used, the geometry can close directly from the tangential condition
\(z\cdot\tau=0\) on each boundary component, without any asymptotics. The present
asymptotic analysis applies to the Gaussian weighted moment/Laplacian formulation
developed above; it should not be used to replace a stronger first-order boundary
condition when that condition is available.
\end{remark}



\section{Finite Weighted Quadrature}

This section gives the precise implication that is actually available. It is a statement
about the weighted Schwarz primitive
\[
  G(z):=F(z,S(z))=\frac{1-e^{-\pi zS(z)}}{\pi z},
\]
not about the ordinary Schwarz function \(S\). For the Gaussian asymptotic problem the
target object is
\[
  M(z):=e^{-\pi zS(z)}.
\]
The use of \(G\) is only a primitive normalization: \(M=1-\pi zG\), so meromorphic
continuation of \(G\) and meromorphic continuation of \(M\) are equivalent away from the
removable point \(z=0\). Meromorphic continuation of \(S\) is neither needed for these
asymptotics nor implied by meromorphicity of \(M\), because recovering \(S\) requires a
single-valued logarithm of \(M\).

We use the following convention for the weighted Cauchy transform:
\begin{equation}\label{eq:weighted-cauchy-K}
  K_f(\zeta):=\int_U \frac{f(z,\bar z)}{\zeta-z}\dd A(z),
  \qquad f(z,\bar z)=e^{-\pi z\bar z}.
\end{equation}
Let \(F(z,\eta)=(1-e^{-\pi z\eta})/(\pi z)\), with removable value \(F(0,\eta)=\eta\), so
that \(\partial_\eta F=e^{-\pi z\eta}\). If \(S\) is the Schwarz function of an analytic
boundary component, then \(G=F(z,S(z))\) on that component.

\begin{proposition}[Cauchy jump for the weighted primitive]\label{prop:cauchy-jump-G}
Let \(\Gamma\) be a real-analytic boundary component and suppose \(G\) is defined in a
collar of \(\Gamma\). Define
\[
  \mathcal C_G(\zeta):=\frac{1}{2i}\int_{\partial U}\frac{G(z)}{\zeta-z}\dd z,
\]
with \(\partial U\) oriented as the boundary of \(U\). Then, for \(\zeta\notin\partial U\),
\begin{equation}\label{eq:C-G-out-in}
  \mathcal C_G(\zeta)=
  \begin{cases}
  K_f(\zeta), & \zeta\in \C\setminus\overline U,\\[0.4em]
  K_f(\zeta)-\pi F(\zeta,\bar\zeta), & \zeta\in U.
  \end{cases}
\end{equation}
Moreover the non-tangential boundary values satisfy
\begin{equation}\label{eq:C-G-jump}
  \mathcal C_G^{\mathrm{out}}-\mathcal C_G^{\mathrm{in}}=\pi G
  \qquad\text{on }\partial U.
\end{equation}
\end{proposition}

\begin{proof}
For \(\zeta\notin\overline U\), Cauchy--Green applied to
\(F(z,\bar z)/(\zeta-z)\) gives
\[
  \int_U \frac{f(z,\bar z)}{\zeta-z}\dd A(z)
  =
  \frac{1}{2i}\int_{\partial U}\frac{F(z,\bar z)}{\zeta-z}\dd z
  =
  \mathcal C_G(\zeta).
\]
For \(\zeta\in U\), use
\(\partial_{\bar z}(1/(\zeta-z))=-\pi\delta_\zeta\). Then
\[
  \frac{1}{2i}\int_{\partial U}\frac{G(z)}{\zeta-z}\dd z
  =
  \int_U \partial_{\bar z}\left(\frac{F(z,\bar z)}{\zeta-z}\right)\dd A(z)
  =
  K_f(\zeta)-\pi F(\zeta,\bar\zeta).
\]
The jump formula follows from the standard Plemelj jump relation with the present
normalization. Equivalently, for \(G\equiv1\) and a positively oriented circle, the inside
value of \(\mathcal C_G\) is \(-\pi\) and the outside value is \(0\), so the jump is
\(\pi G\); the general case follows locally by flattening the analytic arc.
\end{proof}

\begin{theorem}[Finite weighted quadrature]\label{thm:finite-weighted-G}
Assume \(U\) has real-analytic boundary. Suppose the Gaussian weighted area measure has a
finite quadrature identity: there exist points \(a_j\in U\), integers \(m_j\ge0\), and
constants \(c_{j,\ell}\) such that
\begin{equation}\label{eq:finite-weighted-quadrature}
  \int_U h(z)e^{-\pi|z|^2}\dd A(z)
  =
  \sum_{j=1}^N\sum_{\ell=0}^{m_j}c_{j,\ell}h^{(\ell)}(a_j)
\end{equation}
for every holomorphic \(h\) in a neighborhood of \(\overline U\). Then the exterior
weighted Cauchy transform \(K_f\) is rational:
\begin{equation}\label{eq:K-rational}
  K_f(\zeta)
  =
  \sum_{j=1}^N\sum_{\ell=0}^{m_j}
  c_{j,\ell}\,\partial_z^\ell\left(\frac{1}{\zeta-z}\right)\bigg|_{z=a_j},
  \qquad \zeta\in\C\setminus\overline U.
\end{equation}
Consequently, for each boundary component, the weighted Schwarz primitive
\[
  G(z)=\frac{1-e^{-\pi zS(z)}}{\pi z}
\]
extends meromorphically to the \(U\)-side. More precisely, if \(R(\zeta)\) denotes the
rational function on the right-hand side of \eqref{eq:K-rational}, and
\[
  H(\zeta):=K_f(\zeta)-\pi F(\zeta,\bar\zeta)
  \qquad (\zeta\in U)
\]
is the holomorphic interior branch from \eqref{eq:C-G-out-in}, then
\begin{equation}\label{eq:G-meromorphic-formula}
  G_{\mathrm{ext}}(\zeta)=\frac{R(\zeta)-H(\zeta)}{\pi}
\end{equation}
is a meromorphic continuation of \(G\) to the \(U\)-side. Its poles can occur only at the
quadrature nodes \(a_j\), with order controlled by the corresponding \(m_j\).
\end{theorem}

\begin{proof}
Apply \eqref{eq:finite-weighted-quadrature} with
\(h(z)=(\zeta-z)^{-1}\), which is holomorphic near \(\overline U\) for
\(\zeta\in\C\setminus\overline U\). This gives \eqref{eq:K-rational}.

By Proposition~\ref{prop:cauchy-jump-G}, the outside branch of \(\mathcal C_G\) is
\(K_f\), hence agrees with the rational function \(R\). The inside branch is the holomorphic
function \(H=K_f-\pi F\). The jump relation \eqref{eq:C-G-jump} gives, on the boundary,
\[
  R-H=\pi G.
\]
Therefore \((R-H)/\pi\) is a meromorphic function on the \(U\)-side whose boundary values
agree with \(G\). By uniqueness of analytic continuation from the boundary collar, it is
the desired meromorphic continuation. Since \(H\) is holomorphic, the only possible poles
come from \(R\), i.e. from the quadrature nodes and their derivative orders.
\end{proof}

\begin{remark}
This theorem explains exactly what kind of quadrature assumption implies meromorphicity of
\(G\): a finite weighted quadrature identity, equivalently rationality of the Gaussian
weighted Cauchy transform. It still does not imply that \(S\) is meromorphic, because
\[
  e^{-\pi zS(z)}=1-\pi zG(z),
  \qquad
  S(z)=-\frac{1}{\pi z}\Log(1-\pi zG(z)).
\]
Thus meromorphic \(G\) gives a logarithmic, potentially multivalued description of \(S\),
not a meromorphic one unless additional zero/monodromy conditions hold.
\end{remark}

\section{Exact consequences of real analyticity and of a centered annulus}

This section states only what follows from real-analyticity, plus the extra information that
one has a centered annulus inside the continuation region. We use the pure Gaussian object
\[
  J_n^\Gamma=\int_\Gamma z^{n-1}M(z)\dd z,
  \qquad M(z)=e^{-\pi zS(z)}.
\]
All statements below are contour statements for this one-form. They do not assert a true
asymptotic expansion unless a singularity model is explicitly assumed.

\begin{proposition}[Real analyticity alone]\label{prop:real-analytic-alone}
Let \(\Gamma\) be a real-analytic boundary component. Then its Schwarz function \(S\) is
holomorphic in some collar \(N\) of \(\Gamma\), and therefore
\[
  M(z)=e^{-\pi zS(z)}
\]
is holomorphic in that collar. Consequently, for every closed curve \(C\subset N\)
homologous to \(\Gamma\) in \(N\),
\[
  J_n^\Gamma=\int_\Gamma z^{n-1}M(z)\dd z
  =\int_C z^{n-1}M(z)\dd z,
\]
and hence
\[
  |J_n^\Gamma|\le \operatorname{length}(C)\sup_C|M|\,\bigl(\sup_C|z|\bigr)^{n-1}.
\]
Thus real analyticity alone gives exact local deformation identities and exponential
limsup bounds obtained by choosing admissible contours in the analytic collar. It does not
identify a finite set of contributing points and does not give an asymptotic equivalent.
A true Darboux or steepest-descent expansion requires additional information about the
first obstruction to continuing \(M\) beyond the chosen collar.
\end{proposition}

\begin{proposition}[Exact deformation identity]\label{prop:exact-deformation}
Let \(D\subset\C\) be a domain in which \(M\) is holomorphic, and let \(\Gamma\) and
\(C\) be two closed, piecewise \(C^1\) curves in \(D\) that are homologous in \(D\). Then,
for every \(n\ge1\),
\begin{equation}\label{eq:exact-holomorphic-deformation}
  \int_\Gamma z^{n-1}M(z)\dd z
  =
  \int_C z^{n-1}M(z)\dd z.
\end{equation}
Consequently,
\begin{equation}\label{eq:contour-upper-bound}
  \left|\int_\Gamma z^{n-1}M(z)\dd z\right|
  \le
  \operatorname{length}(C)\,\sup_C|M|\,\bigl(\sup_C|z|\bigr)^{n-1}.
\end{equation}
\end{proposition}

\begin{proof}
The one-form \(z^{n-1}M(z)\dd z\) is holomorphic in \(D\). The equality
\eqref{eq:exact-holomorphic-deformation} is Cauchy's theorem for homologous cycles. The
bound \eqref{eq:contour-upper-bound} is the immediate absolute-value estimate on the
right-hand side.
\end{proof}

\begin{corollary}[What a centered annulus gives]\label{cor:centered-annulus-bound}
Assume \(M\) is holomorphic in a region containing a centered annulus
\[
  A_{r_-,r_+}:=\{r_-<|z|<r_+\},
\]
and assume \(\Gamma\) is homologous in that region to every circle
\[
  C_r:=\{|z|=r\},\qquad r_-<r<r_+.
\]
Then, for every fixed \(r\in(r_-,r_+)\),
\begin{equation}\label{eq:annulus-bound-fixed-r}
  J_n^\Gamma
  =
  \int_{C_r} z^{n-1}M(z)\dd z,
  \qquad
  |J_n^\Gamma|
  \le
  2\pi r^n\sup_{|z|=r}|M(z)|.
\end{equation}
In particular,
\begin{equation}\label{eq:annulus-limsup}
  \limsup_{n\to\infty}|J_n^\Gamma|^{1/n}\le r_-.
\end{equation}
\end{corollary}

\begin{proof}
Apply Proposition~\ref{prop:exact-deformation} with \(C=C_r\). The bound follows because
\(\operatorname{length}(C_r)=2\pi r\) and \(|z|^{n-1}=r^{n-1}\) on \(C_r\). Since the
bound holds for each fixed \(r>r_-\), taking \(n\)-th roots and then letting
\(r\downarrow r_-\) gives \eqref{eq:annulus-limsup}.
\end{proof}

\begin{remark}
This is an effective estimate, but it is not a true asymptotic formula. The centered annulus
provides radial comparison contours. It does not identify the first singularities of \(M\),
and therefore it does not determine the leading coefficient or even guarantee a nonzero
leading term.
\end{remark}

\begin{proposition}[Laurent coefficients]\label{prop:laurent-annulus}
Under the hypotheses of Corollary~\ref{cor:centered-annulus-bound}, write the Laurent
expansion of \(M\) in the centered annulus as
\[
  M(z)=\sum_{m\in\mathbb Z} a_m z^m,
  \qquad r_-<|z|<r_+.
\]
Then
\begin{equation}\label{eq:J-laurent-coeff}
  J_n^\Gamma=2\pi i\,a_{-n}.
\end{equation}
Thus the question of asymptotics is exactly the question of asymptotics of the negative
Laurent coefficients of \(M\) in that annulus.
\end{proposition}

\begin{proof}
By Corollary~\ref{cor:centered-annulus-bound}, for any \(r\in(r_-,r_+)\),
\[
  J_n^\Gamma=\int_{|z|=r}z^{n-1}M(z)\dd z.
\]
Substituting the Laurent expansion and integrating term by term gives
\[
  \int_{|z|=r}\sum_{m\in\mathbb Z}a_m z^{m+n-1}\dd z=2\pi i\,a_{-n}.
\]
\end{proof}

\begin{theorem}[When a true Darboux asymptotic follows]\label{thm:exact-darboux-annulus}
Assume the hypotheses of Corollary~\ref{cor:centered-annulus-bound}. Assume further that
\(M\) continues holomorphically through \(|z|=r_-\) except at finitely many points
\(\alpha_1,\dots,\alpha_N\) on \(|z|=r_-\), and that near each \(\alpha_j\)
\begin{equation}\label{eq:finite-singularity-model}
  M(z)=\sum_{\ell=1}^{q_j}b_{j,-\ell}(z-\alpha_j)^{-\ell}+H_j(z),
\end{equation}
where \(H_j\) is holomorphic near \(\alpha_j\). Suppose also that after crossing small
circles around the \(\alpha_j\), the remaining contour can be placed in \(|z|<r_- -\eta\)
for some \(\eta>0\). Then
\begin{equation}\label{eq:exact-pole-asymptotic-annulus}
  J_n^\Gamma
  =
  2\pi i\sum_{j=1}^N\sum_{\ell=1}^{q_j}
  b_{j,-\ell}\binom{n-1}{\ell-1}\alpha_j^{\,n-\ell}
  +O((r_- -\eta)^n).
\end{equation}
In particular, the points contributing to the leading asymptotics are exactly those
singularities among \(\alpha_1,\dots,\alpha_N\) with largest pole order; all lie on the
inner boundary circle \(|z|=r_-\) of the annulus.
\end{theorem}

\begin{proof}
Deform the circle \(C_r\), \(r_-<r<r_+\), inward across the finitely many singularities
\(\alpha_j\), leaving small positively oriented residue circles around them and a remaining
contour in \(|z|<r_- -\eta\). Cauchy's residue theorem gives the sum of residues of
\(z^{n-1}M(z)\dd z\) at the \(\alpha_j\), plus the integral over the remaining contour. The
remaining contour is \(O((r_- -\eta)^n)\) by \eqref{eq:contour-upper-bound}. For a pole term
\(b_{j,-\ell}(z-\alpha_j)^{-\ell}\),
\[
  \operatorname{Res}_{z=\alpha_j}\left(z^{n-1}b_{j,-\ell}(z-\alpha_j)^{-\ell}\dd z\right)
  =b_{j,-\ell}\binom{n-1}{\ell-1}\alpha_j^{n-\ell}.
\]
This proves \eqref{eq:exact-pole-asymptotic-annulus}.
\end{proof}

\begin{remark}[What mere real analyticity does not give]
Real analyticity of the two boundary components gives local collars and therefore bounds of
the form \eqref{eq:contour-upper-bound}. If a centered annulus is contained in the
continuation region, it gives the sharper coefficient statement
\eqref{eq:J-laurent-coeff} and the limsup bound \eqref{eq:annulus-limsup}. It does not give
the finite singularity hypothesis in Theorem~\ref{thm:exact-darboux-annulus}. Without such
additional continuation information, the set of boundary points \(|z|=r_-\) obstructing
continuation of \(M\) can be non-discrete, and no finite Darboux sum is forced.
\end{remark}

\end{document}
